\section{Corrections, conjectures, and further research}
\label{sec:research-program}

\subsection{Two precise corrections to the pinned text}

\paragraph{A consequence is not an equivalence.}
In \texttt{05-signed-kernels.tex}, immediately after
\texttt{signed:eq:rank-conjecture}, the manuscript calls the pair of rank
formulas equivalent to a statement about the dimension of the
Gaussian-supported row space. If $V_w$ is the row space, projection onto
the non-Gaussian columns gives exactly
\[
 \dim\bigl(V_w\cap\QQ^{\{g\text{-columns}\}}\bigr)
 =\rank A_w-\rank A_w^{\neg g}.
\]
Thus the two individual ranks imply the stated dimension, while their
difference alone does not determine either individual rank. Replace
``Equivalently'' by ``Consequently,'' and replace the following sentence
about the equivalence by a sentence about this dimension formula.
Theorem~\ref{newrank:thm:smith} now proves both ranks, so the conjecture can
also be promoted to a theorem when its proof is integrated. The companion
patch makes the narrow logical correction without assuming a particular
chapter-insertion scheme.

\paragraph{Even modified polylogarithms on real arguments.}
In \texttt{03-algebraic.tex}, equation \texttt{bloch:eq:Pm} defines
\[
 P_m(x)=\operatorname{Re}_m\left(
       \sum_{j=0}^{m-1}\frac{2^jB_j}{j!}\log^j|x|\,\Li_{m-j}(x)
       \right),\qquad
 \operatorname{Re}_m=
 \begin{cases}\Re,&m\text{ odd},\\\Im,&m\text{ even}.
 \end{cases}
\]
The later paragraph on higher golden ladders says that the same
coordinates turn a real weight-six ladder into a $\zeta(6)$ multiple.
With this definition, $P_6$ is identically zero on real arguments.
The text itself correctly states this parity fact a little later. In
particular, for the golden arguments in $(0,1)$ every summand inside the
imaginary part is visibly real. A nonzero multiple of $\zeta(6)$ cannot
be the value of that real-argument modified-polylogarithm combination.

The correction is local: remove the weight-six assertion in these
coordinates, retain the raw classical-polylogarithm identity at whatever
status its independent evidence supports, and check odd-weight conversions
directly from the definition. The supplied replacement text avoids
asserting unverified conversions at the other higher weights. No claim
that the raw weight-six ladder is false follows from this observation.

\subsection{Statements that require qualification, not retraction}

The fifth-index theorem does not contradict the baseline theorem, whose
small-index conclusion is restricted to $n\le4$. It disproves the
unrestricted extrapolation of that conclusion. Similarly, the increasing
middle Lerch branch does not contradict a large-derivative asymptotic
offset formula. It shows that its monotonic direction cannot be imposed
on every branch at every small derivative order.

The CM product theorem is a norm statement unless a phase is supplied.
The real-part convention at discriminant $-15$ is mathematically material:
the factor $7/8$ cannot be omitted. The absolute norm alone does not
license a raw signed product at arbitrary discriminants or weights.

The weight-five witness proves an obstruction only for the specified
linear product matrix. No period-independence claim, no impossibility of
using higher-depth relations, and no negative answer to the $S_4$
conjecture follows. This scope was already correctly emphasized by the
baseline; the shorter integral witness makes it easier to check.

\subsection{A concrete next zero-count conjecture}

\begin{conjecture}[Sixth- and seventh-index first derivatives]
\label{conj:n6n7}
The function $\gamma_6'$ has exactly four positive zeros and
$\gamma_7'$ has exactly five positive zeros. In each case every positive
zero is simple.
\end{conjecture}

The numerical candidates are
\begin{align*}
 n=6:\quad&0.8426787120,\ 1.2904536569,\ 1.7994089556,\ 2.2385946523;\\
 n=7:\quad&0.7781418150,\ 0.9499243123,\ 1.4420456198,\\
           &1.9499308927,\ 2.3354109288.
\end{align*}
These statements are \emph{conjectures}, not further certified theorems.
The supplied diagnostic uses 70 decimal digits, an Euler--Maclaurin
evaluation with $M=24$, $R=16$, and 992 test points at each index: 401
logarithmic points spanning $10^{-3}$ to $10^{3.1}$ and 591 linear points
$j/200$, $60\le j\le650$. Root refinement produces the displayed values.
At first derivative order the spectral summands are negative once
$a>e^n$, so this upper tail is analytically free of zeros. Nevertheless,
the finite scan does not certify the unsampled intervals or exclude a
tangency. A proof should use global saturation at a higher derivative
order and descend through every critical point, as in
Theorem~\ref{n5:thm:transition}.

The conjecture is deliberately finite. These two observations do not
support a simple universal formula for the first-derivative count:
the oscillation near $a=1$, the wide range of large roots at higher
derivative orders, and the count plateaus already proved at index five
make extrapolation hazardous.

\subsection{A focused program for the formal relation system}

\paragraph{Construct certificates without elimination.}
The proof of Theorem~\ref{newrank:thm:smith} identifies a unimodular
coordinate system in which the problem breaks into one- and two-variable
blocks. Turn that proof into an explicit symbolic certificate compiler.
For a target vector, it should either return an integer row combination
(or, in even weight, one for twice the target) or return an integral
annihilating witness. The witness basis is already explicit, so membership
testing can be performed with only $w-1$ pairings in odd weight. Tracking
the affine right sides through the same transformations is the remaining
engineering and symbolic-algebra task.

\paragraph{Identify the missing relation for $S_4$.}
The six-entry witness pairs to $7$ with the target. Any proposed
enlargement must include a relation whose coefficient vector does not
annihilate that witness. This provides an exact preliminary test before
computing a much larger matrix. Candidates include distribution,
argument transformations, and elimination through higher depth. The
research question is not merely whether an enlarged matrix has larger
rank, but which analytically justified relation kills this particular
obstruction and with what right side.

\paragraph{Determine the effect of distribution integrally.}
The $2$-torsion in even weight is a property of the stated product
presentation. Does adjoining the appropriate distribution rows remove
some or all of it? What is the Smith form of that enlarged lattice?
This is a concrete question about an explicitly constructible sequence of
integer modules; no unproved transcendence assumption is required.

\paragraph{Generalize the involution method to other levels.}
At level four, conjugation reduces the colors to six and the residual
action becomes the swap $(u,v)\mapsto(v,u)$. At other small levels,
determine whether an integral color decomposition produces similarly
manageable finite-group actions on homogeneous polynomials. Rational
rank formulas alone would miss torsion introduced by nonunimodular
changes of basis, so an integral formulation should be maintained from
the outset.

\subsection{A focused program for CM identities}

\paragraph{Automate the exact arithmetic factor.}
For fixed even weight, generate $F_w$ in exact rational arithmetic and
evaluate its resultant with a certified Hilbert class polynomial.
At weights four and six the arithmetic is concentrated in
$H_D(0)$ and $H_D(1728)$; in higher weights the zeros of the Eisenstein
series create additional algebraic factors. A systematic collection of
these factors would distinguish patterns forced by modular algebra from
patterns inferred from decimal recognition.

\paragraph{Classify phases and genus refinements.}
Extend the phase analysis along the unit arc to every even weight,
including the additional arc zeros of $E_w$. Extend the proved character
products to nonmaximal orders and combine several independent genus
characters to recover every genus product. For higher genus rank the
plus and minus fibers of one character are unions of genera. These are the natural intermediate
objects between a full class norm and an individual CM evaluation.
Every claimed formula should specify the reduced representative,
automorphy factor, and chosen algebraic embedding.

\paragraph{Move beyond full norms.}
For larger class groups, determine individual algebraic factors relative
to a fixed CM period. For the genus ratios, determine when the uniform
degree bound of Theorem~\ref{cmgenus:thm:degree} is attained, and which
proper subfields occur as the weight changes. A rational class norm does not imply that raw sums of
positive-weight CM values are rational or Galois-invariant. Replacing
such raw sums by explicitly normalized modular quantities is the natural
route to well-posed algebraicity questions.

\subsection{A focused program for zero geometry}

\paragraph{Determine saturation thresholds.}
Set
\[
 K_n^{\mathrm{cl}}=\min\{k\ge1:\gamma_n^{(k)}
       \text{ has }n\text{ positive simple zeros}\}.
\]
The exact values are now $K_n^{\mathrm{cl}}=1$ for $1\le n\le4$ and
$K_5^{\mathrm{cl}}=3$. The baseline concentration theorem implies that
the threshold exists for every fixed index. Find a useful effective
bound in $n$, and then determine whether the sequence of thresholds has
any monotonicity or regular asymptotic behavior. The proved parity and
count monotonicity restrict possible answers but do not determine them.

\paragraph{Make the small-parameter theorem effective.}
Theorem~\ref{thm:small-rho} gives a complete count for sufficiently small
positive $\rho$, but not an explicit threshold. Its proof suggests a
constructive route: isolate the algebraic roots of the integer
polynomial, bound derivatives on disjoint disks, and bound the remaining
series uniformly. Quantifying the separation from $\log2$ is the
sensitive step. Transcendence proves nonvanishing, while a useful
numerical threshold needs an effective separation bound.

\paragraph{Certify the fold coordinates.}
The existence and nondegeneracy of the index-three fold are proved.
The diagnostic predicts
\[
 a_*\approx1.04975305088025546,\qquad
 \rho_*\approx0.999990466837738634.
\]
Produce a rational interval box for this pair by applying an interval
Newton method to $(D,\partial_aD)$, with explicit quadrature or
Euler--Maclaurin remainders in both variables. The determinant of the
Jacobian is $-\partial_\rho D\,\partial_a^2D>0$ at the fold, so the
proved nondegeneracy identifies the right certification problem. The
decimals above are not yet such a box.

\paragraph{Describe the entire deformation, not only its endpoints.}
For odd $n$ at first derivative order,
$\partial_\rho(\rho^{-(e^n-a)}D)<0$. Thus each fixed $a$ changes sign
at most once. Determine the topology of the resulting sign domains:
which negative intervals are created, which can merge, and which
bifurcations are possible under the global multiplicity bound. For
$n=3$, the proved fold controls all parameters after its first appearance
in $[4/5,3/2]$, while the possible earlier configuration on the right
remains a separate question.

\paragraph{Treat the simultaneous large-index regime.}
The baseline concentration expansion fixes $n$ and lets $k$ grow.
The extreme roots of $Q_n$ change with $n$, and the small-parameter
Puiseux exponent $1/(n-k)$ exhibits another nonuniform limit. Determine
which double-scaling regimes preserve the known zero approximations,
and where the right asymptotic description must change. A satisfactory
answer would connect the formal spectral polynomials, reciprocal-gamma
root geometry, and the observed threshold sequence.
